% Generated from a `canvaslms quizzes export` of the practice
% exam by mkexempel.py; do not edit by hand.
\begin{otherlanguage}{swedish}
\subsection{Question 1}
\label{ep1:q1}
\emph{Learning objective 1, 1 point, matching.}

Matcha rätt identiferarnamn med rätt beskrivning.

Match each name with a description; the key:
\begin{description}
  \item[{\texttt{i}}] En räknare i en slinga, en lokal variabel
  \item[{\texttt{Djur}}] Ett klassnamn
  \item[{\texttt{ÅLDER}}] En konstant
  \item[{\texttt{läs\_in()}}] Ett funktionsnamn
  \item[{\texttt{filnamn}}] Ett variabelnamn
\end{description}

\subsection{Question 2}
\label{ep1:q2}
\emph{Learning objective 1, 1 point, multi-answer.}

Vilka av följande är bra variabelnamn på en variabel som lagrar ett namn?

The options; every correct option is marked:
\begin{description}
  \item[{A}] NAMN
  \item[{B \textit{(key)}}] namn
  \item[{C}] sträng
  \item[{D}] n
\end{description}


\subsection{Question 3}
\label{ep1:q3}
\emph{Learning objective 1, 1 point, multi-answer.}

För följande funktion, som låter användaren mata in ett värde för temperatur i grader Celsius och räknar om den till grader Fahrenheit, vilka påstående är korrekta?

\begin{minted}{python}
def FIXA():
  a = int(input("Vilken tempratur vill du omvandla? "))
  return 32 + 1.8 * a
\end{minted}

The options; every correct option is marked:
\begin{description}
  \item[{A}] Funktionen är bra namngiven för den beskriver väl vad funktionen gör och följer riktlinjen/standarden för hur en funktion namnges.
  \item[{B}] Variabelnamnet \textquotedbl{}a\textquotedbl{} är bra, för det är tydligt vad det innehåller.
  \item[{C \textit{(key)}}] Ett bättre variabelnamn för att lagra det inmatade värdet hade varit \textquotedbl{}temperatur\textquotedbl{}.
  \item[{D \textit{(key)}}] Ett mer passande funktionsnamn, som också följer riktlinjerna för hur en funktion ska namnges, hade varit: omvandla.
\end{description}


\subsection{Question 4}
\label{ep1:q4}
\emph{Learning objective 2, 10 points, ordering.}

Placera kodraderna i rätt ordning så att utskriften blir:

\begin{minted}{text}
----
 -----
 ------
 -------
\end{minted}

(Kodraderna saknar indentering.)

The lines to order, as the item lists them:
\begin{description}
  \item[{A}] \leavevmode\par
\begin{minted}{python}
for j in range(1,3):
\end{minted}

  \item[{B}] \leavevmode\par
\begin{minted}{python}
print("-"*j, end = "")
\end{minted}

  \item[{C}] \leavevmode\par
\begin{minted}{python}
print("-"*i)
\end{minted}

  \item[{D}] \leavevmode\par
\begin{minted}{python}
for i in range(1,5):
\end{minted}

\end{description}
The key: the order D, A, B, C.


\subsection{Question 5}
\label{ep1:q5}
\emph{Learning objective 2, 10 points, choice.}

Vad skrivs ut av följande kod?

\begin{minted}{python}
hastighet = 50
if hastighet > 70:
  print("snabbt!")
elif hastighet > 100:
  print("väldigt snabbt!")
elif hastighet < 30:
  print("långsamt")
else:
  print("ganska långsamt")
\end{minted}

The options; the key is marked:
\begin{description}
  \item[{A}] snabbt!
  \item[{B}] väldigt snabbt!
  \item[{C}] långsamt
  \item[{D \textit{(key)}}] ganska långsamt
  \item[{E}] Inget av oavstående
\end{description}


\subsection{Question 6}
\label{ep1:q6}
\emph{Learning objective 2, 10 points, choice.}

Titta på följande kod. Hur många gånger kommer hej skrivas ut om vi sätter antal=25 innan while-slingan?

\begin{minted}{python}
while antal <=15:
    print("hej")
    antal +=1
\end{minted}

The options; the key is marked:
\begin{description}
  \item[{A}] 25
  \item[{B}] 24
  \item[{C}] 15
  \item[{D}] 14
  \item[{E \textit{(key)}}] Inga gånger
\end{description}


\subsection{Question 7}
\label{ep1:q7}
\emph{Learning objective 3, 100 points, choice.}

Du vill skriva ut värdena på alla element i listan frukter. Du vill skriva ut de så att de hamnar på samma rad, med mellanslag emellan varje element, utan [] och andra tecken. Vilken av följande kodsnuttar gör det?

The options; the key is marked:
\begin{description}
  \item[{A}] \leavevmode\par
\begin{minted}{python}
frukter = ["mango", "äpple", "päron"]
print(frukter)
\end{minted}

  \item[{B \textit{(key)}}] \leavevmode\par
\begin{minted}{python}
frukter = ["mango", "äpple", "päron"]
for frukt in frukter:
  print(frukt, end = " ")
\end{minted}

  \item[{C}] \leavevmode\par
\begin{minted}{python}
frukter = ["mango", "äpple", "päron"]
for frukt in frukter:
  print(frukt)
\end{minted}

\end{description}


\subsection{Question 8}
\label{ep1:q8}
\emph{Learning objective 3, 100 points, choice.}

Du skulle vilja be användaren att mata in kilopriset på äpplen. Vilken kod är tydligast för användaren?

The options; the key is marked:
\begin{description}
  \item[{A \textit{(key)}}] \leavevmode\par
\begin{minted}{python}
pris = int(input("Vad kostar äpplen per kilo? "))
\end{minted}

  \item[{B}] \leavevmode\par
\begin{minted}{python}
pris = int(input("äpple? "))
\end{minted}

  \item[{C}] \leavevmode\par
\begin{minted}{python}
pris = int(input("Värde? "))
\end{minted}

  \item[{D}] \leavevmode\par
\begin{minted}{python}
pris = int(input("? "))
\end{minted}

\end{description}


\subsection{Question 9}
\label{ep1:q9}
\emph{Learning objective 3, 100 points, multi-answer.}

Du har en klass som hanterar grönsaker och dess priser (se nedan). Om du skulle vilja skriva ut ett objekt vi skapat med (print). Vilka påstående är sanna?

\begin{minted}{python}
class Grönsak:
  
  def __init__(self, namn, pris):
    self.namn = namn
    self.pris = pris
\end{minted}

The options; every correct option is marked:
\begin{description}
  \item[{A \textit{(key)}}] Då behöver vi definera en \_\_str\_\_- metod eller en \_\_repr\_\_metod.
  \item[{B}] Om följande kod läggs till och körs får vi en utskrift \emph{gurka: 23}

\begin{minted}{python}
gurka = Grönsak("gurka", 23)
print(gurka)
\end{minted}

  \item[{C \textit{(key)}}] Om följande kod läggs till och körs får vi en utskrift i stil med:

\begin{minted}{python}
<__main__.Grönsak object at 0x7f4650fd3c10>
gurka = Grönsak("gurka", 23)
print(gurka)
\end{minted}

  \item[{D}] Programmet skulle ge ett felmeddelande \textquotedbl{}AttributeError\textquotedbl{} om följande kod läggs till och körs:

\begin{minted}{python}
gurka = Grönsak("gurka", 23)
print(gurka.namn)
\end{minted}

\end{description}


\subsection{Question 10}
\label{ep1:q10}
\emph{Learning objective 4, 1\,000 points, rich-fill-blank.}

Om jag vill öppna en fil för \emph{läsning}, vad ska det då stå i rutan?

infil = open(\textquotedbl{}textfil.txt\textquotedbl{}, \fbox{\strut\enspace(1)\enspace})

The key for each blank:
\begin{description}
  \item[(1)] \texttt{\textquotedbl{}r\textquotedbl{}}, accepted as the regular expression \texttt{ *\allowbreak{}(\allowbreak{}mode *\allowbreak{}=\allowbreak{} *\allowbreak{})\allowbreak{}?\allowbreak{}(\allowbreak{}\textquotedbl{}r\textquotedbl{}\textbar{}\allowbreak{}\textquotesingle{}r\textquotesingle{})\allowbreak{} *\allowbreak{}(\allowbreak{},\allowbreak{} *\allowbreak{}encoding=\allowbreak{}(\allowbreak{}\textquotedbl{}.\allowbreak{}*\allowbreak{}\textquotedbl{}\textbar{}\allowbreak{}\textquotesingle{}.\allowbreak{}*\allowbreak{}\textquotesingle{})\allowbreak{})\allowbreak{}?\allowbreak{} *\allowbreak{}}
\end{description}

\subsection{Question 11}
\label{ep1:q11}
\emph{Learning objective 4, 1\,000 points, ordering.}

Placera koden i rätt ordning, så att all data läses in från textfilen och varje inläst rad skrivs ut till användaren.

The lines to order, as the item lists them:
\begin{description}
  \item[{A}] \leavevmode\par
\begin{minted}{python}
fil = open("textfil.txt", "r")
\end{minted}

  \item[{B}] \leavevmode\par
\begin{minted}{python}
rad = fil.readline().strip()
\end{minted}

  \item[{C}] \leavevmode\par
\begin{minted}{python}
  rad = fil.readline().strip()
\end{minted}

  \item[{D}] \leavevmode\par
\begin{minted}{python}
  print(rad)
\end{minted}

  \item[{E}] \leavevmode\par
\begin{minted}{python}
while rad != "":
\end{minted}

\end{description}
The key: the order A, B, E, D, C.


\subsection{Question 12}
\label{ep1:q12}
\emph{Learning objective 4, 1\,000 points, choice.}

Du vill skriva över data som du har lagrat i en lista till en fil, som heter \textquotedbl{}resultat.txt\textquotedbl{}. Vilket av påståendena är korrekt?

The options; the key is marked:
\begin{description}
  \item[{A \textit{(key)}}] Då behöver vi först öppna filen \textquotedbl{}resultat.txt\textquotedbl{} för skrivning.
  \item[{B}] Det går inte att göra om det redan finns en fil som heter \textquotedbl{}resultat.txt\textquotedbl{}
  \item[{C}] Då skriver vi:

\textquotedbl{}resultat.txt\textquotedbl{}.write()
  \item[{D}] .write() kan vi använda för att skriva ut integers, vi behöver inte göra om värden i listan till strängar för att kunna skriva ut de på filen.
\end{description}


\subsection{Question 13}
\label{ep1:q13}
\emph{Learning objective 5, 10\,000 points, rich-fill-blank.}

Du vill ta in ett flyttal av användaren som talar om hur många liter grädde det går för att laga en gräddtårta till 12 personer. Vad ska stå i den tomma rutan?

liter\_grädde = \fbox{\strut\enspace(1)\enspace}(input(\textquotedbl{}Hur många liter grädde? \textquotedbl{}))

The key for each blank:
\begin{description}
  \item[(1)] \texttt{float}, accepted as the regular expression \texttt{ *float *}
\end{description}

\subsection{Question 14}
\label{ep1:q14}
\emph{Learning objective 5, 10\,000 points, choice.}

Om du matar in \textquotedbl{}tjugotre\textquotedbl{} som svar på frågan nedan, vilken typ av fel ska du fånga upp med din except-sats? Vilka kodrader är korrekta?

\begin{minted}{python}
try:
  val = int(input("Hur många liter mjölk köpte du?"))
\end{minted}

The options; the key is marked:
\begin{description}
  \item[{A \textit{(key)}}] \leavevmode\par
\begin{minted}{python}
except ValueError:
  print("Inte ett heltal.")
\end{minted}

  \item[{B}] \leavevmode\par
\begin{minted}{python}
except TypeError:
  print("Inte ett heltal.")
\end{minted}

  \item[{C}] \leavevmode\par
\begin{minted}{python}
except IndentationError:
  print("Inte ett heltal.")
\end{minted}

  \item[{D}] \leavevmode\par
\begin{minted}{python}
except SyntaxError:
  print("Inte ett heltal.")
\end{minted}

\end{description}


\subsection{Question 15}
\label{ep1:q15}
\emph{Learning objective 5, 10\,000 points, multi-answer.}

Vad är det som saknas i den här koden om du vill använda det värdet som användaren matar in som svaret på frågan?

\begin{minted}{text}
input("Har du släckt ljusen? ")
\end{minted}

The options; every correct option is marked:
\begin{description}
  \item[{A \textit{(key)}}] Det inmatade värdet måste då lagras i en variabel, t.ex kan du framför satsen skriva:~ svar =
  \item[{B}] input är inte rätt att använda för att ställa en fråga.
  \item[{C}] input tar bara in integers och här vill vi ju lagra en sträng, så det kommer inte fungera.
  \item[{D}] Vi har glömt skriva str() runt input-satsen.
\end{description}


\subsection{Question 16}
\label{ep1:q16}
\emph{Learning objective 6, 100\,000 points, rich-fill-blank.}

Vad ska det stå i de tomma rutorna för att du ska lyckas definiera en klass som heter Katt och ärver från klassen Djur? [Det kommer inga frågor om arv på det riktiga provet.]

class \fbox{\strut\enspace(1)\enspace}(\fbox{\strut\enspace(2)\enspace}):

The key for each blank:
\begin{description}
  \item[(1)] \texttt{Katt}
  \item[(2)] \texttt{Djur}
\end{description}

\subsection{Question 17}
\label{ep1:q17}
\emph{Learning objective 6, 100\,000 points, choice.}

Om du har en klass som representerar en Bil. Vad skulle byt\_ägare kunna vara för något i klassen?

The options; the key is marked:
\begin{description}
  \item[{A}] Ett attribut
  \item[{B \textit{(key)}}] En metod
  \item[{C}] En funktion
  \item[{D}] Ett klassattribut
\end{description}


\subsection{Question 18}
\label{ep1:q18}
\emph{Learning objective 6, 100\,000 points, choice.}

Om du har följande kod, vad saknas för att du ska kunna sortera objekten i godislista på attributet pris?

\begin{minted}{python}
class Godispåse:
  
  def __init__(self, typ, kostar):
    self.sort = typ
    self.pris = int(kostar)
    
godislista = [Godispåse("nappar", 12), Godispåse("lakritsbåtar", 10), Godispåse("choklad", 15)]
\end{minted}

The options; the key is marked:
\begin{description}
  \item[{A \textit{(key)}}] En \_\_lt\_\_-metod behövs definieras i klassen
  \item[{B}] En lista med objekt går aldrig att sortera så vi kan inte lägga till något för att få det att fungera.
  \item[{C}] En \_\_str\_\_-metod
  \item[{D}] Godispåse har inget attribut pris, så det går inte.
\end{description}


\subsection{Question 19}
\label{ep1:q19}
\emph{Learning objective 7, 1\,000\,000 points, multi-answer.}

Du har skrivit en funktion enligt nedan. Vilka av de följande funktionsanropen är då lämpliga, tar emot returvärdet och ger inga felmeddelande?

\begin{minted}{python}
def matbudget(frukost_pris, lunch_pris, middags_pris, antal_dagar = 1):
  summa = frukost_pris + lunch_pris + middags_pris
  summa += 40 #mellanmål och kaffe
  summa = summa * antal_dagar
  return summa
\end{minted}

The options; every correct option is marked:
\begin{description}
  \item[{A \textit{(key)}}] \leavevmode\par
\begin{minted}{python}
kostnad = matbudget(20,70,100)
\end{minted}

  \item[{B \textit{(key)}}] \leavevmode\par
\begin{minted}{python}
kostnad = matbudget(20,70,100,5)
\end{minted}

  \item[{C}] \leavevmode\par
\begin{minted}{python}
kostnad = matbudget(20,70,100,5, mellan_målspris)
\end{minted}

  \item[{D}] \leavevmode\par
\begin{minted}{text}
matbudget(20,70,100)
\end{minted}

\end{description}


\subsection{Question 20}
\label{ep1:q20}
\emph{Learning objective 7, 1\,000\,000 points, multi-answer.}

Varför är det bra att dela upp sitt program i funktioner? Kryssa för alla alternativ som är korrekta.

The options; every correct option is marked:
\begin{description}
  \item[{A \textit{(key)}}] Det gör koden mer lättläst om den är uppdelade så att varje funktion gör en sak.
  \item[{B \textit{(key)}}] Det gör att vi inte behöver upprepa kod flera gånger, utan kan då istället anropa vår funktion flera gånger.
  \item[{C \textit{(key)}}] Det är lättare att felsöka om vi delat upp vårt program i mindre delar.
  \item[{D}] Man måste ha funktioner, annars kan programmet inte köras.
  \item[{E \textit{(key)}}] Det gör att det är lättare att återanvända delar av sin kod i andra program.
\end{description}


\subsection{Question 21}
\label{ep1:q21}
\emph{Learning objective 7, 1\,000\,000 points, choice.}

Vilket av följande program har bäst koduppdelning? Alla program gör precis samma sak, det är ingen skillnad i funktionalitet.

The options; the key is marked:
\begin{description}
  \item[{A \textit{(key)}}] \leavevmode\par
\begin{minted}{python}
def är_heltal(inmatning):
  try:
    int(inmatning)
    return True
  except ValueError:
    print("Inte en siffra försök igen.")
    return False


def meny():
  val = 0
  while val != 2:
    print("Meny\n1. Val 1.\n2. Avsluta.")
    val = input("Vad vill du göra?")
    while not är_heltal(val):
      val = input("Vad vill du göra?")
    val = int(val)
    if val == 1:
      print("Du valde 1.\n")
    elif val ==2:
      print("Programmet avslutas")
    else:
       print("Fel siffra, välj igen.")
    
meny()
\end{minted}

  \item[{B}] \leavevmode\par
\begin{minted}{python}
def är_heltal(inmatning):
  try:
    int(inmatning)
    return True
  except ValueError:
    print("Inte en siffra försök igen.")
    return False

val = 0
while val != 2:
  print("Meny\n1. Val 1.\n2. Avsluta.")
  val = input("Vad vill du göra?")
  while not är_heltal(val):
    val = input("Vad vill du göra?")
  val = int(val)
  if val == 1:
    print("Du valde 1.\n")
  elif val ==2:
    print("Programmet avslutas")
  else:
    print("Fel siffra, välj igen.")
\end{minted}

  \item[{C}] \leavevmode\par
\begin{minted}{python}
val = 0
while val != 2:
  print("Meny\n1. Val 1.\n2. Avsluta.")
  while True:
    val = input("Vad vill du göra?")
    try:
      val = int(val)
      break
    except ValueError:
      print("Inte en siffra försök igen.")
  if val == 1:
    print("Du valde 1.\n")
  elif val ==2:
    print("Programmet avslutas")
  else:
    print("Fel siffra, välj igen.")
\end{minted}

\end{description}

\end{otherlanguage}
